\section{Introduction}
\label{sec:intro}

Kinase inhibitors are among the most intensely investigated compound classes
in drug discovery, with over 80 FDA-approved agents and several hundred in
clinical development.\cite{Roskoski2024}
Because the human kinome is comprised of more than 500 proteins sharing a
highly conserved ATP-binding site,\cite{Manning2002} the design of inhibitors
with controlled target profiles remains a central challenge.
Selectivity is a primary
criterion in lead optimization and a key determinant in the drug approval 
process.\cite{Bhullar2018}

Despite its importance, no single consistent definition of selectivity exists.
The field employs at least four distinct quantitative measures:
the S-score,\cite{Karaman2008} selectivity entropy,\cite{Uitdehaag2011}
Gini coefficient,\cite{Graczyk2010} and ratio-based measures, a family to which
the field assigns the CATDS score.\cite{Klaeger2017}
These definitions are applied interchangeably in the literature without
systematic justification, and their parameter values (the activity threshold
in the S-score, the baseline in the entropy and Gini coefficient, the reference
off-target in ratio-based measures) are chosen by convention rather than
principle.
Standardizing them is a prerequisite for integrating selectivity data across
laboratories and platforms.

Bosc et al.\cite{Bosc2017} showed that these definitions can produce
substantially different compound rankings on the same dataset, and the
KInhibition platform\cite{KInhibition2018} explicitly acknowledged that
conflicting results across metrics are common.
However, the structural sources of this disagreement, such as ``which binding profile
properties determine when definitions agree or conflict, and why?'', have not
been characterized.

The absence of a principled basis for selectivity measurement is not merely
a practical inconvenience. 
Go/no-go decisions in lead optimization are made on the basis of selectivity
assessments whose sensitivity to definitional choice is unknown. The lack of consistency across 
applications of selectivity definitions for different compounds makes 
comparison very challenging. 

Our analyses suggest that existing selectivity definitions are, in some cases,
measuring entirely different qualities of binding profiles, and therefore
cannot be easily converted from one to another. 
In information theory, Shannon\cite{Shannon1948} settled which function should
measure entropy by writing down a few properties a sensible measure must
have (e.g., that it should be maximal when all outcomes are equally
likely) and showing that only one formula satisfies them. In this work, 
we follow this precedent, characterizing how the
competing measures actually behave, stating the properties a good measure
should have, and proposing a satisfactory formula. 

We use four large profiling datasets for kinase inhibitor selectivity
that span three assay technologies in our characterization.
In service of the overarching goal of developing a principled, standardized
formula for defining selectivity, our contributions are as follows:

\begin{itemize}
    \item[(i)] we show that the
S-score, entropy and Gini agree closely with each other on every one of the four
datasets (Spearman $0.72$--$1.00$), while the fold-selectivity ratio $R(k)$
agrees with none of them ($0.14$--$0.62$). We call the first group the
distribution-based family, and $R(k)$ the ratio-based family. The split
follows what each quantifies: how concentrated a compound's activity is across
the whole panel, against how large the gap is to one chosen off-target;

    \item[(ii)] we identify three sources of instability: definitional noise for
compounds with no binding above the activity threshold, ratio-specific instability
driven by near-tied top targets, and distribution-ratio disagreement for
compounds with one dominant target and broad moderate off-target activity;

    \item[(iii)] we propose four motivated properties: (D1) a reliability threshold condition,
(D2) bounded gap sensitivity, (D3) distributional consistency, and (D4) monotonicity under
weak off-target addition. We then find that none of the four definitions
assessed satisfies all four simultaneously;

    \item[(iv)] we propose a candidate measure: a gated, smooth-hinge
effective-target number that satisfies all four properties and recovers its
full-panel ranking from roughly a third of a typical panel.
\end{itemize}

We demonstrate that candidate measure we propose is neither unique as a measure 
satisfying the four properties jointly, nor is it guaranteed to perform better than 
existing measures in all cases. 

